% Options for packages loaded elsewhere
% Options for packages loaded elsewhere
\PassOptionsToPackage{unicode}{hyperref}
\PassOptionsToPackage{hyphens}{url}
\PassOptionsToPackage{dvipsnames,svgnames,x11names}{xcolor}
%
\documentclass[
  letterpaper,
  DIV=11,
  numbers=noendperiod]{scrartcl}
\usepackage{xcolor}
\usepackage{amsmath,amssymb}
\setcounter{secnumdepth}{-\maxdimen} % remove section numbering
\usepackage{iftex}
\ifPDFTeX
  \usepackage[T1]{fontenc}
  \usepackage[utf8]{inputenc}
  \usepackage{textcomp} % provide euro and other symbols
\else % if luatex or xetex
  \usepackage{unicode-math} % this also loads fontspec
  \defaultfontfeatures{Scale=MatchLowercase}
  \defaultfontfeatures[\rmfamily]{Ligatures=TeX,Scale=1}
\fi
\usepackage{lmodern}
\ifPDFTeX\else
  % xetex/luatex font selection
\fi
% Use upquote if available, for straight quotes in verbatim environments
\IfFileExists{upquote.sty}{\usepackage{upquote}}{}
\IfFileExists{microtype.sty}{% use microtype if available
  \usepackage[]{microtype}
  \UseMicrotypeSet[protrusion]{basicmath} % disable protrusion for tt fonts
}{}
\makeatletter
\@ifundefined{KOMAClassName}{% if non-KOMA class
  \IfFileExists{parskip.sty}{%
    \usepackage{parskip}
  }{% else
    \setlength{\parindent}{0pt}
    \setlength{\parskip}{6pt plus 2pt minus 1pt}}
}{% if KOMA class
  \KOMAoptions{parskip=half}}
\makeatother
% Make \paragraph and \subparagraph free-standing
\makeatletter
\ifx\paragraph\undefined\else
  \let\oldparagraph\paragraph
  \renewcommand{\paragraph}{
    \@ifstar
      \xxxParagraphStar
      \xxxParagraphNoStar
  }
  \newcommand{\xxxParagraphStar}[1]{\oldparagraph*{#1}\mbox{}}
  \newcommand{\xxxParagraphNoStar}[1]{\oldparagraph{#1}\mbox{}}
\fi
\ifx\subparagraph\undefined\else
  \let\oldsubparagraph\subparagraph
  \renewcommand{\subparagraph}{
    \@ifstar
      \xxxSubParagraphStar
      \xxxSubParagraphNoStar
  }
  \newcommand{\xxxSubParagraphStar}[1]{\oldsubparagraph*{#1}\mbox{}}
  \newcommand{\xxxSubParagraphNoStar}[1]{\oldsubparagraph{#1}\mbox{}}
\fi
\makeatother


\usepackage{longtable,booktabs,array}
\newcounter{none} % for unnumbered tables
\usepackage{calc} % for calculating minipage widths
% Correct order of tables after \paragraph or \subparagraph
\usepackage{etoolbox}
\makeatletter
\patchcmd\longtable{\par}{\if@noskipsec\mbox{}\fi\par}{}{}
\makeatother
% Allow footnotes in longtable head/foot
\IfFileExists{footnotehyper.sty}{\usepackage{footnotehyper}}{\usepackage{footnote}}
\makesavenoteenv{longtable}
\usepackage{graphicx}
\makeatletter
\newsavebox\pandoc@box
\newcommand*\pandocbounded[1]{% scales image to fit in text height/width
  \sbox\pandoc@box{#1}%
  \Gscale@div\@tempa{\textheight}{\dimexpr\ht\pandoc@box+\dp\pandoc@box\relax}%
  \Gscale@div\@tempb{\linewidth}{\wd\pandoc@box}%
  \ifdim\@tempb\p@<\@tempa\p@\let\@tempa\@tempb\fi% select the smaller of both
  \ifdim\@tempa\p@<\p@\scalebox{\@tempa}{\usebox\pandoc@box}%
  \else\usebox{\pandoc@box}%
  \fi%
}
% Set default figure placement to htbp
\def\fps@figure{htbp}
\makeatother





\setlength{\emergencystretch}{3em} % prevent overfull lines

\providecommand{\tightlist}{%
  \setlength{\itemsep}{0pt}\setlength{\parskip}{0pt}}



 


\KOMAoption{captions}{tableheading}
\makeatletter
\@ifpackageloaded{caption}{}{\usepackage{caption}}
\AtBeginDocument{%
\ifdefined\contentsname
  \renewcommand*\contentsname{Table of contents}
\else
  \newcommand\contentsname{Table of contents}
\fi
\ifdefined\listfigurename
  \renewcommand*\listfigurename{List of Figures}
\else
  \newcommand\listfigurename{List of Figures}
\fi
\ifdefined\listtablename
  \renewcommand*\listtablename{List of Tables}
\else
  \newcommand\listtablename{List of Tables}
\fi
\ifdefined\figurename
  \renewcommand*\figurename{Figure}
\else
  \newcommand\figurename{Figure}
\fi
\ifdefined\tablename
  \renewcommand*\tablename{Table}
\else
  \newcommand\tablename{Table}
\fi
}
\@ifpackageloaded{float}{}{\usepackage{float}}
\floatstyle{ruled}
\@ifundefined{c@chapter}{\newfloat{codelisting}{h}{lop}}{\newfloat{codelisting}{h}{lop}[chapter]}
\floatname{codelisting}{Listing}
\newcommand*\listoflistings{\listof{codelisting}{List of Listings}}
\makeatother
\makeatletter
\makeatother
\makeatletter
\@ifpackageloaded{caption}{}{\usepackage{caption}}
\@ifpackageloaded{subcaption}{}{\usepackage{subcaption}}
\makeatother
\usepackage{bookmark}
\IfFileExists{xurl.sty}{\usepackage{xurl}}{} % add URL line breaks if available
\urlstyle{same}
\hypersetup{
  pdftitle={LING-2801: Neuroscience of Language},
  colorlinks=true,
  linkcolor={blue},
  filecolor={Maroon},
  citecolor={Blue},
  urlcolor={Blue},
  pdfcreator={LaTeX via pandoc}}


\title{LING-2801: Neuroscience of Language}
\author{}
\date{}
\begin{document}
\maketitle


\textbf{Professor}: Heidi Getz,
Ph.D.~\href{mailto:heidi.getz@georgetown.edu}{\nolinkurl{heidi.getz@georgetown.edu}}\\
Linguistics office: Poulton Hall 223 / Neurology office: Building D,
Room 154\\
\textbf{Class schedule}: Tues/Thurs 9:30-10:45\\
\textbf{Office hours}: TBD\\
\textbf{Prerequisites}: None\\
\textbf{Required text}: Brennan, J. R. (2022). \emph{Language and the
brain: a slim guide to neurolinguistics}. Oxford University Press.
(Additional readings will be posted on Canvas.)

\href{https://georgetown.bncollege.com/course-material-listing-page?utm_campaign=storeId=88131_langId=-1_courseData=LING_2801_01_F26&utm_source=wcs&utm_medium=registration_integration}{Link
to purchase text at bookstore}

\section{\texorpdfstring{\textbf{Overview}}{Overview}}\label{overview}

\subsection{\texorpdfstring{\textbf{Course
Description}}{Course Description}}\label{course-description}

This is an introductory course on language and the brain, designed for
students with no background in neuroscience. Its goal is to introduce
students to the major scientific findings in this area as well as the
scientific process that generated this evidence. Topics include the
brain regions and networks involved in language processing in the
healthy brain; how these structures develop in children; and how
language skills are affected by damage to these structures.

\subsection{\texorpdfstring{\textbf{Learning
Goals}}{Learning Goals}}\label{learning-goals}

The course is designed to achieve the following learning outcomes:

\begin{itemize}
\tightlist
\item
  Knowledge of the basic principles, foundational evidence, and current
  research challenges in the field of language and the brain\\
\item
  Understanding of the scientific process through which researchers
  learn about this field\\
\item
  The ability to understand and interpret scientific evidence about
  language in the brain\\
\item
  An understanding of how research on language and the brain relates to
  broader societal issues
\end{itemize}

\section{\texorpdfstring{\textbf{Course
Requirements}}{Course Requirements}}\label{course-requirements}

\emph{More details about all of the following will be provided in
class.}

\subsection{\texorpdfstring{\textbf{Readings}}{Readings}}\label{readings}

There will be two types of readings. Each week, you will have readings
that provide general background on the topic. These will be chapters
from the Brennan text, encyclopedia articles, or similar sources. You
will also read one or two of the original scientific articles that were
cited in the background reading. See \href{?tab=t.78fg1erk9om5}{Schedule
\& Readings} for week-by-week reading assignments.

\subsection{\texorpdfstring{\textbf{Exams:
60\%}}{Exams: 60\%}}\label{exams-60}

There will be four in-person exams. These are designed to assess your
mastery of course content, as well as your ability to map scientific
evidence to theoretical claims.

\subsection{\texorpdfstring{\textbf{Figure Presentations:
15\%}}{Figure Presentations: 15\%}}\label{figure-presentations-15}

You'll be assigned original research papers to read throughout the
semester. These contain figures, tables, and charts that you will learn
to interpret: these illustrate the methods and results of the
researchers' experiments. In class, you will be asked to present at
least two of these tables/charts to your classmates. Your job will be to
describe, in completely concrete terms, how the study worked (if this is
a methods figure) or what the results were (in a results figure). These
explanations must be completely factual, without any opinion.

\subsection{\texorpdfstring{\textbf{Researcher Interview:
4\%}}{Researcher Interview: 4\%}}\label{researcher-interview-4}

For this assignment, you will conduct a structured interview with one of
the predoctoral or postdoctoral scholars in Georgetown's
\href{https://neurolang.georgetown.edu/}{Neuroscience of Language
training program}. The goal of your interview will be to learn about a
specific research question that this person is trying to answer; to
understand the research approach or method that they are trying to use
in order to answer it; and to identify some things that make it
challenging to conduct this type of research. After the interview, you
will briefly describe to your classmates the researcher's main question
and research approach, and you will submit a written reflection.

\subsection{\texorpdfstring{\textbf{Researcher Panel:
1\%}}{Researcher Panel: 1\%}}\label{researcher-panel-1}

At the end of the semester, we will host a panel discussion with
Neuroscience of Language researchers at various stages of their career,
ranging from graduate students to senior professors. The focus of this
panel will be on the future of the field: we'll ask researchers where
they think the field is going and what challenges the field faces. To
prepare for this, I will ask each student to submit some questions for
panelists in advance, and then the class will vote on which questions to
present to the panelists. You will also submit a written reflection
afterwards.

\subsection{\texorpdfstring{\textbf{Study Proposal:
20\%}}{Study Proposal: 20\%}}\label{study-proposal-20}

Over the course of the semester, you will work in small groups to
identify a big-picture question about language and the brain; narrow
this to an answerable research question; and propose an empirical study
that would answer the question. You will present your proposal to the
whole class and lead group discussion of your idea. At the end of the
semester, you will submit a final written paper that includes details of
your study proposal as well as discussion of anticipated results.

\section{\texorpdfstring{\textbf{Course
Policies}}{Course Policies}}\label{course-policies}

\subsection{\texorpdfstring{\textbf{Names}}{Names}}\label{names}

To help me learn your name, please consider using
\href{https://www.speakpipe.com/HeidiGetz}{this link} to record yourself
saying your name the way you'd like me to pronounce it.

\subsection{\texorpdfstring{\textbf{Attendance}}{Attendance}}\label{attendance}

Attendance is mandatory. Unexplained absences will affect your grade.
That said:

\begin{itemize}
\tightlist
\item
  If you know in advance that you won't be able to make it to a
  particular class due to a religious holiday, sports conflict, etc., no
  problem---just let me know ahead of time.\\
\item
  If you're sick and/or contagious, please DO NOT come to class. Email
  me to let me know you are not well enough to come. I will not ask any
  follow-up questions, and I will never require you to provide a
  doctor's note.
\end{itemize}

There is no Zoom option for this course.

\subsection{\texorpdfstring{\textbf{Use of
AI}}{Use of AI}}\label{use-of-ai}

I think that generative AI (ChatGPT, etc.) can be a great tool and that
students should learn to use it. However, this course is not the place
for you to do that. This course is also not one in which generative AI
would do a particularly good job. In any case, I expect you to complete
all of your work without the aid of generative AI, unless I specifically
say otherwise. \textbf{In this course, presenting AI-generated work
(including text and ideas) as your own constitutes a violation of
academic integrity}. If I suspect that AI was used for an assignment, I
will refer the submission to the Honor Council for investigation. This
can have enormous consequences for you, including a permanent notation
on your transcript or even suspension or dismissal from Georgetown.

It is ok to use basic copy-editing/spell-checking tools like the ones
provided by Grammarly, Microsoft Word, and Google Docs.

\subsection{\texorpdfstring{\textbf{Accommodations}}{Accommodations}}\label{accommodations}

If you might require an accommodation in this course, you may contact
the \href{https://academicsupport.georgetown.edu/}{Academic Resource
Center (ARC)}. Available accommodations include notetaking assistance,
interpreting/CART services, extra time on assignments, or any other
accommodation that the ARC deems ``reasonable'' based on a conversation
with you and documentation from the appropriate people (e.g.~your
doctor). To arrange accommodations (or just to explore this option),
email \href{mailto:arc@georgetown.edu}{\nolinkurl{arc@georgetown.edu}}
or visit the ARC in-person (Leavey Center, Suite 335, which is
accessible from the elevators next to the student credit union, around
the corner from the bookstore). There's more information about this
process on \href{https://academicsupport.georgetown.edu/disability/}{the
ARC's website}.

Recording course lectures: One possible ARC accommodation is access to
special note-taking software. This software allows students to
audio-record class, generate a transcription of that recording, and
store it alongside other course materials (e.g., lecture slides). When
students are approved for this accommodation, they must sign a document
agreeing that recordings are for their personal use only, and they are
required to delete all recordings at the end of the semester.

\subsection{\texorpdfstring{\textbf{Communication}}{Communication}}\label{communication}

To promote the highest degree of education possible, we ask each student
to respect the opinions, thoughts, and identity of other students and be
courteous in the way that you choose to express yourself. We should all
be respectful and considerate of all opinions.

\subsection{\texorpdfstring{\textbf{Religious
holidays}}{Religious holidays}}\label{religious-holidays}

Georgetown University promotes respect for all religions. Any student
who is unable to attend classes or to participate in any examination,
presentation, or assignment on a given day because of the observance of
a major religious holiday or related travel shall be excused and
provided with the opportunity to make up, without unreasonable burden,
any work that has been missed for this reason and shall not in any other
way be penalized for the absence or rescheduled work. Students will
remain responsible for all assigned work. Students should notify
professors in writing at the beginning of the semester of religious
observances that conflict with their classes. The Office of the Provost,
in consultation with Campus Ministry and the Registrar, will publish,
before classes begin for a given term, a list of major religious
holidays likely to affect Georgetown students. This list can be found on
the Campus Ministry website,
\href{http://campusministry.georgetown.edu/}{http://campusministry.georgetown.edu}.
The Provost and the Main Campus Executive Faculty encourage faculty to
accommodate students whose bona fide religious observances in other ways
impede normal participation in a course. Students who cannot be
accommodated should discuss the matter with an advising dean.

\section{\texorpdfstring{\textbf{Resources}}{Resources}}\label{resources}

\subsection{\texorpdfstring{\textbf{Sexual
Misconduct}}{Sexual Misconduct}}\label{sexual-misconduct}

I am committed to supporting survivors of sexual misconduct, including
relationship violence, sexual harassment, and sexual assault.
\textbf{However, university policy also requires me to report any
disclosures about sexual misconduct to the Title IX Coordinator, whose
role is to coordinate the University's response to sexual misconduct.}
If I report an incident to the Title IX Coordinator, the coordinator
would then reach out to you to provide support, resources, and the
option to meet. You would not be required to meet with the coordinator.
More information about reporting options and resources is available
\href{https://sexualassault.georgetown.edu/\#_ga=2.84763163.302405512.1597672771-114879328.1580830596}{here}.

As an alternative to speaking with me or another mandated reporter, you
can reach out for support and assistance to these other fully
confidential professional resources:

\begin{itemize}
\tightlist
\item
  Confidential email address for Sexual Assault Response and Prevention
  (SARP):
  \href{mailto:sarp@georgetown.edu}{\nolinkurl{sarp@georgetown.edu}}\\
\item
  Jen Schweer, MA, LPC, Associate Director of Health Education Services
  for Sexual Assault Response and Prevention (202) 687-0323 \textbar{}
  \href{mailto:jls242@georgetown.edu}{\nolinkurl{jls242@georgetown.edu}}\\
\item
  Brittany Egan, MSW, LICSW, CCTP, Staff Clinician and Sexual Assault
  Services Specialist, Health Education Services (202) 687-0323
  \textbar{}
  \href{mailto:be200@georgetown.edu}{\nolinkurl{be200@georgetown.edu}}\\
\item
  Counseling and Psychiatric Services (CAPS), (202) 687-6985
\end{itemize}

\subsection{\texorpdfstring{\textbf{Office of the Student Ombuds
(OSO)}}{Office of the Student Ombuds (OSO)}}\label{office-of-the-student-ombuds-oso}

Confidential \textbar{} Independent \textbar{} Impartial \textbar{}
Informal

The Office of the Student Ombuds (OSO) serves all undergraduate and
graduate students, including SCS and BGE, on the main campus. Consider
contacting the Student Ombuds when you want to talk to a caring
professional about a University-related issue but don't know where to
turn. The OSO is a confidential and safe space that is independent of
formal university organizations or structures where students can discuss
their concerns, share their experiences, ask questions and explore their
options. The student ombuds can help you problem-solve, identify your
goals, and empower you to think through ways to navigate complex
situations. Some reasons for you to visit the office may be to address
academic concerns, clarify administrative policies, discuss
interpersonal conflicts, seek coaching, mediation or facilitation to
handle a sensitive situation, advise you on the process to file a formal
complaint if you are experiencing bias, harassment, bullying or other
forms of intimidation, identify other appropriate campus resources, and
allow you to safely express your frustrations and concerns.

Request an in-person or zoom appointment with the Student Ombuds by
writing studentombuds@georgetown.edu or calling 202-784-1081. The OSO is
located in Room 207 of the Reiss Building (across from Arrupe Hall).
Find more information at
\href{http://studentombuds.georgetown.edu/}{http://studentombuds.georgetown.edu}.

\subsection{\texorpdfstring{\textbf{Learning
Resources}}{Learning Resources}}\label{learning-resources}

Georgetown offers a host of
\href{http://scs.georgetown.edu/academic-affairs/resources/}{learning
resources} to its students. Two that you might find particularly helpful
in this course are the
\href{https://writingcenter.georgetown.edu/}{Writing Center} and
\href{http://guides.library.georgetown.edu/refworks}{Refworks}.

\begin{itemize}
\tightlist
\item
  \href{https://writingcenter.georgetown.edu/}{The Writing Center}
  offers peer tutoring by trained graduate and undergraduate students
  who can assist you at any point in the writing process. They help at
  any stage of your writing process, from brainstorming to revision.
  Tutors can offer advice on thesis development, use of evidence,
  organization, flow, sentence structure, grammar, and more. The Writing
  Center will not proofread or edit papers; rather, they will help to
  improve your proofreading and editing skills to become a better
  writer. Appointments can be booked online through their website.\\
\item
  \href{http://guides.library.georgetown.edu/refworks}{Refworks} is an
  online research management tool that aids in organizing, storing, and
  presenting citation sources for papers and projects.
\end{itemize}




\end{document}
